% Regulamin Stowarzyszenia „Iławski Instytut Idei” – wierne przepisanie dokumentu złożonego w ewidencji
% (pisownia i numeracja punktów jak w oryginale, w tym brak punktu 24).
% Budowanie: npm run regulamin  →  src/static/dokumenty/regulamin.pdf
\documentclass[11pt,a4paper]{article}
\usepackage{fontspec}
\usepackage{polyglossia}
\setdefaultlanguage{polish}
\usepackage[margin=2.5cm]{geometry}
\usepackage{enumitem}
\usepackage{lastpage}
\usepackage{fancyhdr}
\usepackage[hidelinks]{hyperref}

\hypersetup{
  pdftitle={Regulamin Stowarzyszenia „Iławski Instytut Idei”},
  pdfauthor={Stowarzyszenie Iławski Instytut Idei},
  pdflang={pl},
}

\pagestyle{fancy}
\fancyhf{}
\renewcommand{\headrulewidth}{0pt}
\fancyfoot[L]{\small Regulamin Stowarzyszenia „Iławski Instytut Idei”}
\fancyfoot[R]{\small \thepage/\pageref{LastPage}}

\setlength{\parindent}{0pt}
\setlength{\emergencystretch}{2em}
\setlist{itemsep=.3em, topsep=.3em}
\setlist[enumerate,1]{label=\textbf{\arabic*.}, leftmargin=2.2em, itemsep=.8em}
\setlist[enumerate,2]{label=\arabic*., leftmargin=1.8em}

\begin{document}

\begin{center}
  {\Large\bfseries REGULAMIN STOWARZYSZENIA „IŁAWSKI INSTYTUT IDEI”}
\end{center}
\vspace{1em}

\begin{enumerate}
  \item Stowarzyszenie zwykłe nosi nazwę „IŁAWSKI INSTYTUT IDEI”. W~dalszych postanowieniach regulaminu zwane jest Stowarzyszeniem.

  \item Terenem działania Stowarzyszenia jest Polska. Dla realizacji celów statutowych Stowarzyszenie może prowadzić działania poza granicami kraju, na terenie innych państw, obszarów, z~poszanowaniem tamtejszego prawa.

  \item Siedzibą Stowarzyszenia jest miejscowość Iława.

  \item Stowarzyszenie jest zawiązane na czas nieokreślony. Działa na podstawie niniejszego regulaminu oraz przepisów Ustawy z~dnia 7~kwietnia 1989~r. Prawo o~stowarzyszeniach.

  \item Działalność Stowarzyszenia oparta jest na pracy społecznej jej członków.

  \item Cele Stowarzyszenia to:
  \begin{enumerate}
    \item Pomoc mieszkańcom Iławy i~okolic, w~rozwoju opartym na ekologi, nauce, nowoczesnych technologiach, aktywnej przedsiębiorczości i~praktycznej edukacji.
    \item Pomoc mieszkańcom Polski w~organizacji społeczności wspierających członków we wzajemnym rozwoju opartym na ekologi, nauce, nowoczesnych technologiach, aktywnej przedsiębiorczości i~praktycznej edukacji.
    \item Realizacja działań wymienionych w~par.~7
  \end{enumerate}

  \item Stowarzyszenie realizuje swoje cele poprzez:
  \begin{enumerate}
    \item Działalność wspomagającą rozwój wspólnot i~społeczności lokalnych;
    \item Działalność w~zakresie ekologii i~ochrony zwierząt oraz ochrony dziedzictwa przyrodniczego;
    \item Działalność wspomagającą rozwój gospodarczy, w~tym rozwój przedsiębiorczości;
    \item Działalność wspomagającą rozwój techniki, wynalazczości i~innowacyjności oraz rozpowszechnianie i~wdrażanie nowych rozwiązań technicznych w~praktyce gospodarczej;
    \item Działalność wspomagającą rozwój turystyki i~krajoznawstwa;
    \item Działalność w~zakresie ochrony i~promocji zdrowia;
    \item Działalności na rzecz dzieci i~młodzieży, w~tym wypoczynku dzieci i~młodzieży;
    \item Działalność wspomagającą rozwój nauki, szkolnictwa wyższego, edukacji, oświaty i~wychowania;
    \item Udzielania nieodpłatnego poradnictwa obywatelskiego.
  \end{enumerate}

  \item Szczegółowa realizacja celów Stowarzyszenia będzie odbywać się poprzez:
  \begin{enumerate}
    \item organizowanie warsztatów twórczych i~edukacyjnych;
    \item organizowanie konferencji, seminariów, szkoleń i~wykładów;
    \item organizowanie imprez kulturalnych i~rozrywkowych;
    \item organizowanie i~rozwój przedsięwzięć edukacyjnych;
    \item działalność informacyjną i~wydawniczą;
    \item organizowanie spotkań;
    \item aktywizację społeczności;
    \item tworzenie nowych koncepcji i~polityk w~dziedzinie współpracy w~zakresie edukacji, technologii i~ekologii;
    \item oddziaływanie na instytucje publiczne i~samorządowe, poprzez doradztwo, konsultacje i~edukację opracowywanie ekspertyz.
  \end{enumerate}

  \item Członkiem stowarzyszenia może być osoba fizyczna posiadająca pełną zdolność do czynności prawnych i~niepozbawiona praw publicznych, będąca obywatelem polskim lub cudzoziemcem.

  \item Przyjęcia nowych członków dokonuje Zebranie Członków uchwałą podjętą w~ciągu 30~dni od daty złożenia deklaracji zawierającej rekomendację dwóch członków.

  \item Członek ma prawo:
  \begin{enumerate}
    \item głosu w~sprawach Stowarzyszenia;
    \item udziału we wszystkich formach działalności Stowarzyszenia;
    \item biernego i~czynnego uczestniczenia w~wyborach do władz Stowarzyszenia;
    \item do otrzymywania wszystkich informacji dotyczących działalności Stowarzyszenia;
  \end{enumerate}

  \item Członek obowiązany jest do:
  \begin{enumerate}
    \item przestrzegania postanowień Regulaminu Stowarzyszenia i~uchwał władz Stowarzyszenia;
    \item aktywnego udziału w~realizacji celów statutowych Stowarzyszenia oraz czynnego udziału w~pracach prowadzonych przez Stowarzyszenie;
    \item regularnego uiszczania składek członkowskich;
  \end{enumerate}

  \item Członkostwo w~Stowarzyszeniu ustaje na skutek:
  \begin{enumerate}
    \item dobrowolnej rezygnacji pisemnej z~przynależności do Stowarzyszenia złożonej na ręce Przedstawiciela
    \item wykluczenia przez Zebrania Członków: a.~z~powodu nieusprawiedliwionego zalegania z~opłatą składek członkowskich lub innych zobowiązań, przez okres przekraczający sześć miesięcy, b.~z~powodu rażącego naruszenia zasad Regulaminu, nieprzestrzegania uchwał władz Stowarzyszenia, c.~ze względu na brak przejawów aktywnej działalności na rzecz Stowarzyszenia;
    \item utraty praw obywatelskich na mocy prawomocnego wyroku sądu, śmierci członka lub utraty osobowości prawnej przez członka.
  \end{enumerate}

  \item Osoba wykluczona lub skreślona ma prawo wniesienia odwołania do Zebrania Członków w~terminie 14~dni od daty doręczenia uchwały. Zebranie Członków rozpatruje odwołanie w~czasie najbliższych obrad, a~jego decyzja jest ostateczna.

  \item Władzami Stowarzyszenia są:
  \begin{enumerate}
    \item Zebranie Członków,
    \item Przedstawiciel.
  \end{enumerate}

  \item Decyzje Zebrania Członków podejmowane są w~głosowaniu jawnym. Zebranie może zdecydować o~przeprowadzeniu głosowania tajnego.

  \item
  \begin{enumerate}
    \item Uchwały Zebrania Członków Stowarzyszenia podejmowane są zwykłą większością głosów przy obecności co najmniej połowy osób uprawnionych do głosowania.
    \item W~sprawach określonych w~par.~19 pkt.~6--10 wymagana jest obecność wszystkich członków stowarzyszenia.
  \end{enumerate}

  \item Zebranie Członków jest najwyższą władzą Stowarzyszenia, które podejmuje decyzje we wszystkich sprawach należących do zakresu działania Stowarzyszenia, o~ile regulamin nie stanowi inaczej.

  \item Do kompetencji Walnego Zebrania należy m.\,in:
  \begin{enumerate}
    \item ustalanie kierunków działania i~rozwoju
    \item wybór i~odwołanie Przedstawiciela
    \item przyjmowanie i~odwoływanie członków
    \item uchwalanie zmian regulaminu
    \item podejmowanie decyzji o~rozwiązaniu stowarzyszenia
    \item podejmowanie decyzji o~nabyciu oraz zbyciu nieruchomości lub prawa użytkowania wieczystego
    \item podejmowanie decyzji o~ustanowieniu ograniczonego prawa rzeczowego
    \item podejmowanie decyzji o~zawarciu umowy kredytu albo pożyczki
    \item podejmowanie decyzji o~przejęciu długu, uznaniu długu, zwolnieniu z~długu, przystąpieniu do długu, zawarciu umowy poręczenia lub zawarciu innej podobnej umowy
    \item podejmowanie decyzji o~zaciąganiu innych zobowiązań przekraczających wartość 1000~zł
  \end{enumerate}

  \item Zebranie Członków zwołuje przedstawiciel lub 2/3 członków Stowarzyszenia. Zebraniem Członków kieruje Przewodniczący Zebrania wybrany na początku zebrania, zwykłą większością głosów.

  \item Stowarzyszenie jest reprezentowane przez Przedstawiciela wybieranego przez Zebranie Członków na 3~letnią kadencję. Przedstawiciel może zostać odwołany przez Zebranie Członków.

  \item Do zakresu działania Przedstawiciela należy:
  \begin{enumerate}
    \item reprezentowanie stowarzyszenia na zewnątrz
    \item zarządzanie majątkiem stowarzyszenia
    \item zwoływanie zebrania członków
  \end{enumerate}

  \item Podjęcie przez przedstawiciela decyzji przekraczającej uprawnienia zwykłego zarządu wymaga zgody wszystkich członków stowarzyszenia. Są to w~szczególności:
  \begin{enumerate}
    \item nabycie oraz zbycie nieruchomości lub prawa użytkowania wieczystego
    \item ustanowienie ograniczonego prawa rzeczowego
    \item zawarcie umowy kredytu albo pożyczki
    \item przejęcie długu, uznanie długu, zwolnieniu z~długu, przystąpieniu do długu, zawarciu umowy poręczenia lub zawarciu innej podobnej umowy
    \item zaciąganie innych zobowiązań przekraczających wartość 1000~zł
  \end{enumerate}

  % W oryginale po punkcie 23 następuje punkt 25.
  \setcounter{enumi}{24}
  \item Środki na działalność stowarzyszeniach pochodzą z:
  \begin{enumerate}
    \item składek członkowskich
    \item dotacji
    \item darowizn
    \item zbiórek publicznych
    \item spadków, zapisów,
    \item dochodów z~majątku stowarzyszenia
  \end{enumerate}

  \item Decyzję w~sprawie zmiany regulaminu oraz rozwiązania stowarzyszenia podejmuje Zebranie Członków zwykłą większością głosów przy obecności co najmniej 2/3 członków uprawnionych do głosowania. Podejmując uchwałę o~rozwiązaniu Stowarzyszenia Zebranie Członków określa przeznaczenie pozostałych środków finansowych Stowarzyszenia.

  \item W~sprawach nieuregulowanych niniejszym Regulaminem zastosowanie mają przepisy ustawy -- Prawo o~stowarzyszeniach.
\end{enumerate}

\end{document}
